%% Triangle des changements de base, sur l'exemple du nombre 59.
%% 59 = (0011 1011)_2 = (3B)_16 ; les blocs de 4 bits donnent les symboles hexadécimaux.
\documentclass[tikz,border=4pt]{standalone}
\usepackage{liaisons}
\begin{document}
\begin{tikzpicture}[x=1cm,y=1cm,
  boite/.style={draw, line width=1.2pt, rounded corners=3pt, fill=white,
                minimum width=2.6cm, minimum height=1cm, align=center, font=\footnotesize},
  fl/.style={-{Stealth[length=5pt]}, line width=1pt},
  etiq/.style={font=\scriptsize, align=center, text width=2cm, inner sep=1pt,
               execute at begin node={\hyphenpenalty=10000\exhyphenpenalty=10000}}]

%% ================== les trois écritures ===============================
\node[boite, draw=solideB] (dix)  at (3.4,3.0) {Base 10\\$59_{(10)}$};
\node[boite, draw=solideE] (deux) at (0,0.5)   {Base 2\\$(0011\,1011)_2$};
\node[boite, draw=solideC] (hexa) at (6.8,0.5) {Base 16\\$(3\mathrm{B})_{16}$};

%% ================== descentes : divisions successives =================
\draw[fl, solideE] (2.05,2.5) -- (0.2,1.05);
\node[etiq, text=solideE, anchor=south east] at (1.05,1.9)
     {restes des divisions successives par 2};
\draw[fl, solideC] (4.75,2.5) -- (6.6,1.05);
\node[etiq, text=solideC, anchor=south west] at (5.75,1.9)
     {restes des divisions successives par 16};

%% ================== remontées : somme des poids =======================
\draw[fl, solideB] (1.15,1.05) -- (2.75,2.5);
\draw[fl, solideB] (5.65,1.05) -- (4.05,2.5);
\node[etiq, text=solideB, fill=white, text width=1.7cm] at (3.4,1.75)
     {somme des poids $\times$ chiffre};

%% ================== équivalence binaire / hexadécimal =================
\draw[{Stealth[length=5pt]}-{Stealth[length=5pt]}, line width=1pt, solideF] (1.4,0.5) -- (5.4,0.5);
\node[font=\scriptsize, text=solideF, anchor=south, inner sep=2pt] at (3.4,0.56) {blocs de 4 bits};

%% ================== tableaux de poids =================================
\foreach \dx/\bits/\somme in {-0.9/{0,0,1,1}/{$2+1=3$}, 4.3/{1,0,1,1}/{$8+2+1=11=\mathrm{B}$}} {
  \begin{scope}[shift={(\dx,-1.6)}]
    \foreach \p [count=\i from 0] in {8,4,2,1} {
      \node[font=\scriptsize] at (\i*0.45+0.225,0.62) {\p};
      \draw[black!55, line width=0.6pt] (\i*0.45,0) rectangle (\i*0.45+0.45,0.45); }
    \foreach \b [count=\i from 0] in \bits {
      \node[font=\scriptsize] at (\i*0.45+0.225,0.225) {\b}; }
    \node[font=\scriptsize, anchor=west] at (1.95,0.225) {\somme};
  \end{scope} }
\end{tikzpicture}
\end{document}
